\documentclass[a4paper,10pt]{article}
\usepackage[T1]{fontenc}
\usepackage[utf8]{inputenc}
\usepackage{lmodern,textcomp}
\usepackage[margin=2.2cm]{geometry}
\usepackage{amsmath,amssymb,siunitx,booktabs,graphicx,tabularx,array,enumitem,xcolor,fancyhdr}
\usepackage[colorlinks=true,linkcolor=blue!50!black,urlcolor=blue!50!black]{hyperref}
\definecolor{draft}{RGB}{135,70,0}
\setlength{\parindent}{0pt}
\setlength{\parskip}{4pt}
\raggedright
\setlength{\headheight}{14pt}
\renewcommand{\arraystretch}{1.2}
\renewcommand{\tabularxcolumn}[1]{>{\raggedright\arraybackslash}p{#1}}
\setlist{nosep,leftmargin=*}
\newcommand{\todo}[1]{\par\smallskip{\color{draft}\textbf{To complete:} #1}\par\smallskip}
\newcommand{\blank}{\textcolor{draft}{\textit{[fill in]}}}
\newcommand{\file}[1]{\nolinkurl{#1}}
\newcommand{\slot}[2]{\begin{center}\fbox{\begin{minipage}[c][25mm][c]{0.91\linewidth}\centering
\textbf{Figure to add}\\#1\\{\small\file{#2}}\end{minipage}}\end{center}}
\pagestyle{fancy}\fancyhf{}
\fancyhead[L]{34870 Electroacoustics}
\fancyhead[R]{Lab C --- report skeleton}
\fancyfoot[L]{Group 10 \quad / \quad Measurements pending}
\fancyfoot[R]{\thepage}
\graphicspath{{../figures/}{Lab C/figures/}}
\title{34870 Electroacoustics\\Lab C --- Microphone calibration}
\author{Group 10: Louis Andrianne, Sophie Kimura, Mads Rudolph}
\date{Laboratory date: 29 September 2026}
\begin{document}
\begin{center}
{\Large\bfseries Lab C --- Microphone calibration}\\[5pt]
34870 Electroacoustics\\[4pt]
Group 10: Louis Andrianne, Sophie Kimura, Mads Rudolph\\[4pt]
Laboratory date: Tuesday 29 September 2026, 10:30--12:30 (building 354, room 026)
\end{center}\medskip
\thispagestyle{fancy}
\textbf{Working draft.} This skeleton follows the official E26 lab brief. Replace the
coloured prompts and figure boxes with your group's evidence. Given settings and theory
are not measured results. Keep raw data unchanged and record which files were analysed.

\section{Purpose and measurement model}
The exercise combines single-frequency absolute sensitivity calibration with broadband
electrostatic-actuator measurements of two condenser microphones. The corrected responses
are used to estimate resonance frequency and damping and to fit a condenser-microphone
model in LTspice. The metrological interpretation includes repeatability, reproducibility
and uncertainty.

\par\smallskip
\begin{tabularx}{\linewidth}{@{}lXX@{}}\toprule
Microphone & Serial number in brief & Confirmed capsule / notes \\\midrule
B\&K 4133 & 591628 & \blank \\
B\&K 4134 & 1534527 & \blank \\\bottomrule
\end{tabularx}\par\smallskip
The actuator provides a response shape; the absolute scale is set using a calibrator or
pistonphone measurement. Let
\begin{equation}
H_{21}(f)=\frac{V_2(f)}{V_1(f)},\qquad
H_{pv}(f)=\frac{H_{21}(f)}{H_{21,\mathrm{ref}}(f)}.
\end{equation}
$H_{21,\mathrm{ref}}$ characterizes the acquisition/amplifier chain. The corrected actuator
ratio $H_{pv}$ is not itself an absolute sensitivity in volts per pascal.

\section{Equipment and preparation (Part 0)}
\par\smallskip
\begin{tabularx}{\linewidth}{@{}lX@{}}\toprule
Equipment & Identity / actual setting \\\midrule
PC / MATLAB / NI USB-4431 & \blank \\
B\&K Nexus / microphone preamplifier & \blank \\
GRAS 14AA / B\&K actuator & \blank \\
GRAS 42AG calibrators (three units) & Serial numbers to record in Part 2 \\
Pistonphone, if available & Type and serial: \blank \\
Room temperature / static pressure & \blank \\\bottomrule
\end{tabularx}\par\smallskip
\textbf{Part 0 check:} connect AO0 to AI0 and AI1, then run
\file{test_meas_mag_spec2}. Record whether acquisition works and identify any changes.

\textbf{Handling constraints from the brief.} Never touch the diaphragm. Keep the
actuator supply off during positioning or capsule changes and on only while measuring.
Disconnect the microphone cable from the Nexus before changing capsules. Switch the
Nexus on before reconnecting the preamplifier, then verify its settings.
\todo{Confirm the equipment identities and successful
initial connection test. Add a labelled setup photo/sketch for each measurement configuration.}

\newpage
\section{Part 1: Acquisition-system reference}
Split AO0 to AI0 and the Nexus input using the specified cables; connect the Nexus
output to AI1. This measures the chain needed to correct subsequent voltage ratios.

\par\smallskip
\begin{tabularx}{\linewidth}{@{}lXX@{}}\toprule
Nexus setting & Specified value & Actual / checked \\\midrule
Channel & 1 & \blank \\
Output setting & \SI{10}{\milli\volt\per\pascal} & \blank \\
Input sensitivity & \SI{10}{\milli\volt\per\pascal} & \blank \\
Resulting gain $G_{\mathrm{Nex}}$ & 1 & \blank \\
High-pass / low-pass & \SI{0.1}{\hertz} / \SI{100}{\kilo\hertz} & \blank \\
Polarization voltage & \SI{200}{\volt} & \blank \\\bottomrule
\end{tabularx}\par\smallskip

\par\smallskip
\begin{tabularx}{\linewidth}{@{}lXlX@{}}\toprule
Parameter & Brief value & Parameter & Brief value \\\midrule
$f_1$ / $f_2$ & 20 / \SI{60000}{\hertz} & $A_{\max}$ & 1 \\
Tones/octave & 6 & $N_{\mathrm{av}}$ & 4 \\
Boost corner $f_b$ & \SI{1}{\hertz} & Resolution & \SI{1}{\hertz} \\\bottomrule
\end{tabularx}\par\smallskip
\textbf{Bandwidth check before interpreting plots.} The brief requests an upper
frequency of \SI{60}{\kilo\hertz}, while the repository copy of
\file{meas_mag_spec2.m} sets the sample rate to \SI{96}{\kilo\hertz}. The Nyquist
limit is therefore \SI{48}{\kilo\hertz}. Check the actual routine and valid generated
tones with the TA; do not treat frequencies at or above Nyquist as independently measured
response points. Record the usable acquisition band and any changed settings.

Calculate and save the complex \file{H_21_ref = specn(:,2)./specn(:,1)}.
\par\smallskip
\begin{tabularx}{\linewidth}{@{}lXXX@{}}\toprule
Frequency & $|H_{21,\mathrm{ref}}|$ & Phase / degrees & Extraction / interpolation \\\midrule
\SI{250}{\hertz} & \blank & \blank & \blank \\
\SI{1000}{\hertz} & \blank & \blank & \blank \\\bottomrule
\end{tabularx}\par\smallskip
\slot{Reference spectra plus system-transfer magnitude and phase.}{labC_system_reference.png}
\todo{State the saved reference file, actual frequency band and interpolation method.
Explain why dividing by a complex reference corrects both magnitude and phase.}

\newpage
\section{Part 2: Single-frequency sensitivity calibration}
Recheck the Nexus settings, slowly seat each capsule in the correct adaptor and leave
the assembly undisturbed. Use the frequency and level on each device's sticker or
certificate, rather than assuming a nominal level.
\begin{align}
p_{\mathrm{rms}}&=p_{\mathrm{ref}}10^{L_p/20},& p_{\mathrm{ref}}&=\SI{20}{\micro\pascal},\\
M(f_c)&=\frac{\operatorname{Calibrate}(f_c,L_p,G_{\mathrm{Nex}})}
{|H_{21,\mathrm{ref}}(f_c)|},&
L_M&=20\log_{10}\!\left(\frac{M}{\SI{1}{\volt\per\pascal}}\right).
\end{align}
The course function already applies its supplied amplifier gain; do not correct that
gain a second time. The wrapper logs corrected $M$ in V/Pa.

\par\smallskip
\begin{tabularx}{\linewidth}{@{}lXXXX@{}}\toprule
Device & Serial no. & $f_c$/Hz & Certified $L_p$/dB & Corrections / uncertainty \\\midrule
42AG unit 1 & \blank & \blank & \blank & \blank \\
42AG unit 2 & \blank & \blank & \blank & \blank \\
42AG unit 3 & \blank & \blank & \blank & \blank \\
Pistonphone$^*$ & \blank & \blank & \blank & \blank \\\bottomrule
\end{tabularx}\par\smallskip
{\small $^*$Subject to availability. Record any pressure correction required by its certificate.}

\subsection{Corrected sensitivity results}
Calibrate both microphones with every available device. For a repeatability estimate,
repeat each pairing under the same conditions; record whether the capsule was removed
and reseated between runs. Add rows or use a separate raw-data table as needed.

\par\smallskip
\begin{tabularx}{\linewidth}{@{}llXXXX@{}}\toprule
Mic. & Device & $n$ & Mean $M$ & Std. dev. $s_M$ & $L_M$ / dB \\
& & & mV/Pa & mV/Pa & re 1 V/Pa \\\midrule
4133 & 42AG 1 & \blank & \blank & \blank & \blank \\
4133 & 42AG 2 & \blank & \blank & \blank & \blank \\
4133 & 42AG 3 & \blank & \blank & \blank & \blank \\
4133 & Pistonphone & \blank & \blank & \blank & \blank \\
4134 & 42AG 1 & \blank & \blank & \blank & \blank \\
4134 & 42AG 2 & \blank & \blank & \blank & \blank \\
4134 & 42AG 3 & \blank & \blank & \blank & \blank \\
4134 & Pistonphone & \blank & \blank & \blank & \blank \\\bottomrule
\end{tabularx}\par\smallskip
\todo{Populate from \file{labC_calibration_log.csv}. Group repeats by microphone,
device, frequency and level. Do not silently combine results at different frequencies.}
\todo{Compare repeated runs and different devices. Separate repeatability from changes
in measurement conditions; explain why none of these estimates establishes an exact true value.}

\newpage
\section{Uncertainty of the single-frequency calibration}
State the measurand, microphone, calibration frequency and environmental conditions.
For an underlying voltage model write
\begin{equation}
M=\frac{V}{G\,|H_{21,\mathrm{ref}}|\,p_{\mathrm{ref}}10^{L_p/20}}.
\end{equation}
For example, $\partial M/\partial L_p=-(\ln 10/20)M$ when $L_p$ is expressed in dB.
For general inputs $x_i$ and sensitivity coefficients $c_i=\partial M/\partial x_i$,
\begin{equation}
u_c^2(M)=\sum_i c_i^2u^2(x_i)+2\sum_{i<j}c_ic_j\operatorname{cov}(x_i,x_j).
\end{equation}
Omit covariance terms only with a justified independence assumption. Avoid counting the
same DAQ/amplifier contribution again if it is already included in the measured reference.

\par\smallskip
\begin{tabularx}{\linewidth}{@{}lXXX@{}}\toprule
Input / contribution & Evidence / distribution & Standard uncertainty & Contribution to $u(M)$ \\\midrule
Repeated calibration & Same-condition repeats & $s_M/\sqrt{n}$ for mean & \blank \\
Calibrator level & Certificate; convert stated $U/k$ & \blank & \blank \\
Reference correction & Repeat/reference data & \blank & \blank \\
Residual gain / voltage & Specification or verification & \blank & \blank \\
Seating / environment & Experimental estimate & \blank & \blank \\\bottomrule
\end{tabularx}\par\smallskip
\todo{Fill separate budgets if the two capsules or calibration devices require them.
Use actual certificates: no generic instrument accuracy is a substitute for evidence.
Distinguish spread of individual readings from uncertainty of their mean.}

\subsection{Report the selected calibration result}
\par\smallskip
\begin{tabularx}{\linewidth}{@{}lXX@{}}\toprule
Quantity & B\&K 4133 & B\&K 4134 \\\midrule
Selected calibration device & \blank & \blank \\
Calibration frequency $f_c$ & \blank & \blank \\
$M(f_c)$ / mV Pa$^{-1}$ & \blank & \blank \\
$u_c(M)$ / mV Pa$^{-1}$ & \blank & \blank \\
Expanded $U=k u_c$, stated $k$ & \blank & \blank \\\bottomrule
\end{tabularx}\par\smallskip
\todo{Justify the calibration result selected to scale the actuator measurement. Discuss
its traceability and explain the chosen coverage factor. If evidence is missing, identify
the missing contribution rather than inventing a numerical uncertainty.}

\newpage
\section{Part 3: Actuator response}
\subsection{Setup and noise floor (Parts 3a--3b)}
Split AO0 to AI0 and the GRAS 14AA input. The actuator excites the diaphragm; the
microphone/preamplifier feeds the Nexus, whose output goes to AI1. Keep the room quiet
and apply the handling precautions on page 1.

With the actuator supply \emph{off}, use Part 1's settings except
$f_b=\SI{200}{\hertz}$. Save noise measurements for $N_{\mathrm{av}}=4$ and 64.
With the supply \emph{on}, acquire the microphone response with the specified settings,
then repeat for the other capsule. Verify the valid bandwidth as in Part 1.
\par\smallskip
\begin{tabularx}{\linewidth}{@{}lXXX@{}}\toprule
Condition & Capsule & Averages & Actual file / repeat \\\midrule
Actuator off & \blank & 4 & \blank \\
Actuator off & \blank & 64 & \blank \\
Actuator on & 4133 & \blank & \blank \\
Actuator on & 4134 & \blank & \blank \\\bottomrule
\end{tabularx}\par\smallskip
\slot{Noise at both averaging settings and actuator-on signal, same voltage convention.}{labC_noise_floor.png}
For coherent averaging of uncorrelated noise, the expected amplitude-level change is
$-10\log_{10}(64/4)\simeq-\SI{12.0}{\decibel}$. This is a prediction, not a measured reduction.
\todo{Quantify the observed noise change and signal clearance over a stated band.
Identify hum or correlated disturbances that do not follow the ideal averaging law.}

\subsection{Correct and scale the responses (Part 3c)}
Use complex division $H_{pv}=H_{21}/H_{21,\mathrm{ref}}$ on matching frequency grids.
For a selected sensitivity magnitude $M(f_c)$, the broadband magnitude is
\begin{equation}
|M(f)|=M(f_c)\frac{|H_{pv}(f)|}{|H_{pv}(f_c)|}.
\end{equation}
Use the corrected phase with a documented reference/sign convention. A single-frequency
magnitude calibration does not independently establish absolute phase.
\todo{State the selected $f_c$, calibration value and interpolation. The current
\file{process_labC.m} averages all calibration rows and scales a fitted low-frequency
shape; verify or adapt this before claiming exact scaling at your selected $f_c$.}
\todo{Compare the response shapes and identify the free-field and pressure-field
microphones from the evidence. Explain the relation to the scattering correction in Lab B.}

\newpage
\section{Resonance frequency, Q and LTspice model}
\subsection{Estimate resonance and damping}
For a second-order low-pass approximation,
\begin{equation}
\frac{M(f)}{M_0}=\frac{1}{1-(f/f_s)^2+\mathrm{j}(f/f_s)/Q}.
\end{equation}
$f_s$ occurs at a $90^\circ$ phase lag relative to the low-frequency plateau, and
$Q=|M(f_s)|/|M_0|$. $Q$ is a linear ratio, not a dB value. Use the course slide
procedure first; a separate complex fit can be a cross-check.
\par\smallskip
\begin{tabularx}{\linewidth}{@{}lXXXX@{}}\toprule
Capsule & Phase-reference band & $f_s$/Hz & $Q$ & Fit cross-check \\\midrule
4133 & \blank & \blank & \blank & \blank \\
4134 & \blank & \blank & \blank & \blank \\\bottomrule
\end{tabularx}\par\smallskip
\todo{Show the phase crossing and amplitude ratio. State interpolation, fitting range
and uncertainty; avoid a reference band already affected by damping or low-frequency roll-off.}

\subsection{Given model parameters}
\par\smallskip
\begin{tabularx}{\linewidth}{@{}lXlX@{}}\toprule
Parameter & Brief value & Parameter & Brief value \\\midrule
Gap $x_0$ & \SI{20.77}{\micro\metre} & Polarization $E$ & \SI{200}{\volt} \\
Effective diameter & \SI{8.95}{\milli\metre} & Diaphragm mass $M_{MD}$ & \SI{1.5}{\milli\gram} \\
Compliance $C_{MD}$ & \SI{0.02}{\milli\metre\per\newton} & Damping $R_{MD}$ & \SI{0}{\newton\second\per\metre} \\
Back volume $V_B$ & \SI{126.4}{\milli\metre\cubed} & Air properties $\rho,c$ & State assumptions \\\bottomrule
\end{tabularx}\par\smallskip
The repository's model uses
\begin{align}
S_D&=\pi(D_{\mathrm{eff}}/2)^2,& C_{AB}&=\frac{V_B}{\rho c^2},&
C_{MT}&=\left(C_{MD}^{-1}+S_D^2/C_{AB}\right)^{-1},\\
M_{MT}&=\frac{1}{(2\pi f_s)^2 C_{MT}},&
M_{AS}&=\frac{M_{MT}-M_{MD}}{S_D^2}-M_{A1},&
R_{AS}&=\frac{1}{Q S_D^2}\sqrt{\frac{M_{MT}}{C_{MT}}}.
\end{align}
Here $M_{A1}$ is the front acoustic mass used by the chosen circuit; document that
term and keep its definition consistent when extracting the backplate mass.
\textbf{No free-field scattering source:} omit the $G_{pb}$ generator implementing
$T(s)$ because the actuator measurement has no incident sound field.
\todo{Insert the circuit with labelled acoustic, mechanical and electrical domains.
The existing \file{LabC_CondenserMics.asc} is a starting point, not a fitted measurement result.}

\newpage
\subsection{Initial and fitted backplate parameters}
\par\smallskip
\begin{tabularx}{\linewidth}{@{}lXXXX@{}}\toprule
Capsule & Initial $M_{AS}$ & Fitted $M_{AS}$ & Initial $R_{AS}$ & Fitted $R_{AS}$ \\
& kg/m$^4$ & kg/m$^4$ & Pa s/m$^3$ & Pa s/m$^3$ \\\midrule
4133 & \blank & \blank & \blank & \blank \\
4134 & \blank & \blank & \blank & \blank \\\bottomrule
\end{tabularx}\par\smallskip
\todo{Explain any adjustment from the values inferred from $f_s$ and $Q$. Check signs,
units and physical plausibility before interpreting an inferred mass or resistance.}

\subsection{Measurement versus simulation}
\slot{Both capsules: calibrated magnitude and corrected phase, overlaid with LTspice output.}{labC_ltspice_comparison.png}
\todo{Generate/export actual LTspice traces and overlay them with the measurement.
The existing \file{labC_responses_vs_model.png} is a second-order model comparison;
it does not by itself document an LTspice run. State any amplitude normalization and
any $180^\circ$ phase adjustment caused by the ground convention.}

\par\smallskip
\begin{tabularx}{\linewidth}{@{}lXX@{}}\toprule
Check & 4133 & 4134 \\\midrule
Frequency band compared & \blank & \blank \\
Magnitude agreement / residual & \blank & \blank \\
Phase agreement / residual & \blank & \blank \\
Main remaining discrepancy & \blank & \blank \\\bottomrule
\end{tabularx}\par\smallskip
The brief's appendix notes that this approximate model may predict a low-frequency
sensitivity below the nominal \SI{12.5}{\milli\volt\per\pascal}. Its effective acoustic
and capacitor areas are not exactly the same. Do not force the model to the nominal
sensitivity or capacitance while concealing the resulting model limitation.

\section{Discussion and conclusion}
\todo{Summarize the calibrated sensitivities with uncertainties, the evidence for each
microphone's field type, the measured $f_s$ and $Q$, and the validity of the fitted model.
Distinguish observed results, model predictions and unresolved limitations. Explain what
the actuator can establish and why an acoustic calibration is still needed.}

\section*{Sources and processing record}
\small
DTU 34870, \emph{Lab C -- Microphone calibration -- E26},
\file{34870_Lab_C_ActuatorCalibration_E2026.pdf}, pp. 1--6; Lectures 4 and 6 for the
resonance/Q method. Source: course vault / DTU Learn.

Repository: \file{Lab C/matlab/measure_labC.m}, \file{calibrate_labC.m},
\file{process_labC.m}, and \file{Lab C/ltspice/LabC_CondenserMics.asc}.

\textbf{Reference/response repeats, calibration rows, processing and circuit versions:} \blank
\end{document}
